\chapter{看懂分子模型}

\begin{kaichang}
翻开一页药物说明书、一篇科普文章，或者这本书后面的章节，你很可能会撞见这样的东西：几个字母之间拉着短线；一串像锯齿的折线，折线尽头蹲着一两个字母；几根小棍穿着彩色小球；还有一团鼓鼓囊囊、互相挤在一起的“泡泡”。它们声称画的是同一个东西——一个分子。

同一个分子，为什么要画出四五种完全不同的样子？哪一种才是分子“真实的长相”？

答案是：都不是，也都是。这一篇附录教你认读五种最常见的分子“画像”：路易斯结构、骨架式、楔线式、球棍模型、空间填充模型。读完你会发现，每张图都是一场交易：它保留一些信息，同时故意扔掉另一些。学会看懂它们，关键不是背画法，而是看懂每笔交易的条件——它想让你看见什么？它为了不让你分心，藏起了什么？
\end{kaichang}

\section{先说清楚：图不是分子本身}

五种画像背后有一条共同原则：\textbf{它们都是人为约定的表示法}。真实原子没有颜色，小球涂成黑色、红色、蓝色只是为了区分元素；化学键不是小棍也不是短线，它是电子在原子之间的一种排布方式（第 17、19 章）；分子也不是静态照片，它们在不停地振动、旋转、弯折。每种模型都像一张地图：地图把立体的城市压扁在纸上、删掉每条小巷，才换来“一眼看清主干道”的便利。看懂分子模型的第一步，是先问一句：这张图拿什么东西，换了什么东西？

下面五种模型，从“最省墨”到“最占地方”，一个比一个接近分子的外形，也一个比一个看不清内部结构。我们就按这个顺序认识它们。

\section{路易斯结构：电子的账本}

路易斯结构（第 19 章讲过它的来历和局限）是五种模型里最“会计”的一种。它要回答的问题只有一个：\textbf{价电子在哪儿}。

读法只有三条：

\begin{itemize}[itemsep=2pt]
\item 元素符号代表原子核加上内层电子；
\item 两个原子之间的一条短线，代表一对被共享的电子，也就是一根共价键；两条线是双键，三条线是三键；
\item 画在原子旁边的圆点是没有参与成键的“孤对电子”——这对电子常常是反应里的关键角色，第 32 章的路易斯酸碱就靠它定义。
\end{itemize}

以水为例：氧在中间，左右各拉一条线连向氢，氧的上下再各点一对点。两秒钟你就读出：两根 \chem{O-H} 单键、两对孤对电子。再看二氧化碳：碳居中，左右各两条线连向氧，每个氧上还有两对点——碳氧之间是双键，所有原子周围的电子数都对得上账。路易斯结构的本事就在这里：它把每个价电子的去向记成一笔明账，让你能检查“账平不平”。

线的数目本身就是信息。你呼吸的氧气 \chem{O_2}，两个氧之间是两条线（双键）；空气里最多的氮气 \chem{N_2}，两个氮之间是三条线（三键）。三键把两个氮锁得极牢，拆开它们要花很大的力气——这就是为什么氮气占了空气的约 $78\%$，却几乎不参与任何反应，植物想利用它还得靠根瘤菌或化肥厂帮忙（第 86 章）。一张小小的路易斯结构，已经把“氮气为什么懒惰”的答案写在纸上了。

\textbf{保留了什么}：原子之间谁连谁（连接方式）、有几重键、每个原子周围还剩几对孤对电子、电荷落在谁头上。这些信息足以让你推断很多反应会怎样发生。

\textbf{牺牲了什么}：几乎牺牲了全部几何信息。路易斯结构把水画成一条直线或直角，可真实水分子两根 \chem{O-H} 键的夹角约为 $104.5^{\circ}$——第 20 章讲过，正是这个“弯”决定了水的极性，进而决定了水的几乎全部特殊性质（第 23 章）。此外，圆点暗示电子是静止的小珠子，而第 19 章强调过，电子其实是弥散的云；遇到苯这类电子“匀”在整个环上的分子，一张路易斯结构根本画不下，得靠几幅共振结构轮流值班（第 19、41 章）。

还有一招常考你：带电的离子怎么画？约定是把整个结构括在方括号里，电荷标在右上角。比如水的“亲戚”氢氧根 \chem{OH^-}：氧和氢之间一条线，氧上三对孤对电子，外面套一个方括号、右上角一个负号——这表示整个离子多出来一个电子。画碳酸根 \chem{CO_3^{2-}} 时你会发现一幅画不完：双键可以在三个氧之间“轮流坐庄”，这逼着化学家发明共振结构（第 19 章），也预告了路易斯账本的局限。

什么时候用它：数电子、查电荷、写反应机理的时候。它是账本，不是照片。

\section{骨架式：有机化学家的速写}

翻开任何一本有机化学书或一篇新药新闻，满纸都是锯齿线。这就是骨架式（也叫键线式），第 36 章认识碳链时已经见过它。碳氢化合物里的碳和氢多得写不过来，骨架式就是为这种分子发明的速写。

读法也只有三条，但第一次见的人几乎都会误读：

\begin{itemize}[itemsep=2pt]
\item 折线的每一个\textbf{端点和折点}都是一个碳原子——碳字根本不写出来；
\item 连在碳上的氢也全省略。你默认碳总是伸出四根键（第 36 章），没有被其他原子占用的键，全都连着氢；
\item 碳、氢以外的原子（氧、氮、氯……）必须明明白白写出符号，连在它们身上的氢通常也写出来，比如 $-\chem{OH}$。
\end{itemize}

以乙醇为例：骨架式就是一折两段的小短线，末端写一个 $\chem{OH}$。两段线、一个折点、三个字母，就把 \chem{C_2H_6O} 的全部连接细节交代清楚了——两个碳、碳上连着五个氢、末端那个碳还牵着一个羟基。再比如第 40 章的乙酸（醋的酸味来源）：在乙醇的速写旁多画一个碳氧双键，两个分子的“性格差异”在图上一目了然。

骨架式还有两个常考你的细节。一是环：一个封闭的六边形就是一个六元环。如果六边形里再画一个圆圈或三根相间的短线，那是苯环（第 41 章），六个碳共享一团离域电子——这时千万别去数“哪条边是双键”，圆圈就是告诉你“别数了，电子是匀的”。二是双线：两段平行线表示双键，比如羧基里那根碳氧双键。试着读一读阿司匹林（一种常见止痛药）的骨架式：一个苯环，环上相邻两个位置各伸出一小段——一段连着羧基，一段连着一个小小的酯基。学会这两招，药物说明书上大半的分子图你都能读出结构来。

\textbf{保留了什么}：碳骨架的连接方式（直链、支链还是环）和官能团的位置——而第 37 章说过，官能团正是有机分子的“反应接口”。画成锯齿还不只是省墨：绕着单键，碳链可以旋转（第 36 章），锯齿形大致反映了碳链伸展时一种常见的姿态。

\textbf{牺牲了什么}：全部碳氢细节（你得自己在脑子里把氢补回来）、原子的真实大小、准确的键角（纸上约 $120^{\circ}$ 的折角不等于真实的约 $109.5^{\circ}$），以及所有三维信息——骨架式把分子压平在纸面上，读它的人必须自己记得分子其实是立体的。手性这种纯三维的信息在普通骨架式里完全丢失，要请下一种模型出场。

什么时候用它：追踪大分子碳骨架的变化、认官能团、快速比较两个分子是不是同一种连接（第 36 章的同分异构）的时候。

\section{楔线式：把纸面折出三维}

骨架式骗你说分子是平的，楔线式负责拆穿这个谎言。它在骨架式的基础上增加两种特殊的线：

\begin{itemize}[itemsep=2pt]
\item 一条\textbf{实心三角楔}：这根键从纸面里伸出来，指向你；
\item 一条\textbf{虚线楔}（一排由短到长的横线）：这根键缩到纸面后面去，远离你；
\item 普通细线：这根键躺在纸面所在的平面里。
\end{itemize}

它最常登场的场合是第 42 章的手性。一个碳连着四个不同的基团时，这四种基团在空中有两种互为镜像、无法重合的排法，就像你的左手和右手。普通平面式里，这两种分子画得一模一样；换成楔线式——一个基团画实心楔、一个画虚线楔、两个画普通线——左右手立刻分得清清楚楚。比如乳酸：一种排法出现在你运动后酸胀的肌肉里，另一种是它的镜像。你身体里的酶只认其中一种，因为酶本身也是有手性的“口袋”（第 42、49 章）。

读楔线式有个实用技巧：先找“焦点原子”——那几根实心楔和虚线楔交汇的点，通常就是手性中心；再分别确认朝你的、背你的、躺在纸面上的各是什么基团。有些教材还会把楔线画得由窄到宽，宽的一头离你近。记住纸面只是一个参照系，把整幅图在脑子里旋转一下，前后关系不变——这正是在训练你第 42 章需要的“三维眼力”。

\textbf{保留了什么}：关键原子上各基团的前后指向，也就是立体化学信息。这在药物里是要命的信息：第 42 章讲过，一对镜像分子的药效可以天差地别，一种治病，另一种可能无效甚至有害。

\textbf{牺牲了什么}：它只标出“焦点”附近的三维关系，分子的其余部分依然是压平的；楔线只告诉你相对的前后，不给精确的键角和距离。还有一个更隐蔽的坑：绕单键的旋转时刻在进行（第 36 章），楔线式冻结的只是无数姿态中的一个瞬间，别把那张快照当成分子的全部生活。

什么时候用它：讨论手性、比较镜像分子、追踪反应过程中基团朝向变化的时候。

\section{球棍模型：看得见的“骨架”}

从这里开始，我们从“纸面符号”走向“实物形象”。球棍模型大概是最著名的分子形象：彩色小球代表原子，小棍代表共价键；球上按特定角度钻好孔，小棍插进去，自然张开正确的键角——碳的四面体角约 $109.5^{\circ}$ 就是这样体现的。你几乎能在每张化学实验室的照片里看到它。

颜色是约定俗成的（不同厂商略有出入）：碳黑、氢白、氧红、氮蓝、硫黄、氯绿。请把这句话抄在脑子里：\textbf{这是约定，不是原子的颜色}。原子比可见光的波长还小几百倍，“颜色”对它根本没有意义。

\begin{qiaoliang}
甲烷那四根键为什么张成约 $109.5^{\circ}$？这不是审美，是能量账。碳周围有四对成键电子，电子带负电、彼此排斥，四对电子离得越开，整个体系的总能量越低。在平面上，四对电子最远只能张成 $90^{\circ}$ 的十字；把它们支出纸面、指向正四面体的四个顶点，两两夹角就变成约 $109.5^{\circ}$——这才是“最宽敞”的排法。第 20 章用这个思路（电子对互斥）解释了水为什么弯、氨为什么是三角锥：分子形状的图纸，写在电子的排斥账里。
\end{qiaoliang}

\textbf{保留了什么}：连接方式、键角、整体几何构型。把模型拿在手里一转，你能直观体会第 20 章那句话——分子的形状决定它会做什么。

\textbf{牺牲了什么}：首先是比例。小棍又硬又长，暗示相邻原子之间隔着空旷的距离；其实成键原子的电子云紧紧贴在一起，真实分子里既没有那根“棍”，也没有棍所占的空间（第 17 章专门拆穿过“键是小棍”的想象）。其次是电子：孤对电子、电子云的弥散，在球棍模型里彻底隐身。

什么时候用它：入门阶段讲形状、讲键角、讲“这个分子能不能正好卡进那个口袋”的时候。

\section{空间填充模型：分子真正的“身材”}

球棍模型把原子缩小、把键拉长，空间填充模型反其道而行：把每个原子按真实的相对大小膨胀成球，球与球直接挤在一起，棍子消失。你看到的是一团凹凸起伏的“泡泡”——这是五种模型里最接近分子“外形”的一种。第 48 章里折叠成团的蛋白质、第 49 章里酶的活性口袋、第 66 章里 DNA 双螺旋的沟槽，画的多半就是它。

做个数量级估算，你就明白它为什么更“诚实”。典型共价键的键长在 $10^{-10}$ 米量级，比如一根 \chem{C-H} 键约 \ud{1.1\times 10^{-10}}{m}；而碳原子本身的半径也有这个量级的一半以上。也就是说，两个成键原子的中心距离，比两个原子半径之和还要小——两个球必然深深嵌进彼此，根本分不出缝隙。球棍模型里“两个小球隔棍相望”的画面，在真实尺度下不存在；空间填充模型那团互相吞掉的泡泡，才接近分子占据空间的真实情况。

\textbf{保留了什么}：分子占据的空间、表面的形状和起伏。药物分子能不能滑进酶的口袋、底物能不能通过细胞膜的通道（第 51 章）、水分子怎样贴在蛋白质表面——这类“接触”问题，只有空间填充模型答得上。

\textbf{牺牲了什么}：你看不见键了，也看不见埋在内部的原子。一个蛋白质的空间填充图就是一团“菜花”，什么连接方式都读不出来；键角、键级、孤对电子全部失踪。它诚实于外形，却对外形之下的一切守口如瓶。

什么时候用它：关心分子之间的“接触”——识别、结合、堵塞、通过——的时候。

\section{实战演练：同一个分子，五种读法}

光看定义不够，我们拿一个真实分子走一遍全套流程。就选丙氨酸——组成你身体里蛋白质的二十种氨基酸之一（第 48、69 章）。它的“档案”：一个中心碳，连着四样东西：一个氨基（$-\chem{NH_2}$）、一个羧基（$-\chem{COOH}$）、一个甲基（$-\chem{CH_3}$）、一个氢。

\textbf{第一步，路易斯结构}。把每个原子摆出来，逐对清点价电子：中心碳四根单键，氮上一对孤对电子，羧基里一根碳氧双键、每个氧上两对孤对电子。这步做完，账平了，你也知道了：氮上那对孤对电子使它有本事接受一个质子——这正是氨基酸既能当酸又能当碱的原因（第 32 章）。

\textbf{第二步，骨架式}。碳字和碳上的氢全部省略：一条小折线代表甲基连着中心碳，中心碳再伸出两笔，分别写上 $\chem{NH_2}$ 和 $\chem{COOH}$。三秒钟画完，碳骨架和三个官能团的位置一目了然。

\textbf{第三步，楔线式}。中心碳连着四个不同的基团——这是个手性中心。把氨基画成实心楔（朝向你）、氢画成虚线楔（藏到纸后）、甲基和羧基画成普通线，得到一种丙氨酸；把楔的方向对调，得到它的镜像。你身体里的蛋白质，几乎只用其中一种（第 42 章）。

\textbf{第四步，球棍模型}。黑球是碳、白球是氢、蓝球是氮、红球是氧，小棍按约 $109.5^{\circ}$ 张开。拿在手里一转，你能看到氨基和羧基在空间里的相对方位——这是平面图给不了的直觉。

\textbf{第五步，空间填充模型}。所有原子按真实大小膨胀、挤在一起：小棍消失，只剩一团带凹凸的“泡泡”，氨基那里微微凸起一块蓝色。这团泡泡的形状，决定了丙氨酸能不能嵌进某个酶的口袋（第 49 章）。

五种读法，五种问题。你以后在任何书里遇到丙氨酸，看到的都是这五张画像中的一张——认出是哪张，你就知道作者想让你看什么。

\section{五种模型的对账表}

\begin{table}[htbp]
\centering
\begin{tabular}{llll}
\toprule
模型 & 回答的问题 & 保留了什么 & 牺牲了什么 \\
\midrule
路易斯结构 & 价电子在哪儿 & 连接方式、键级、孤对电子 & 三维形状、电子的弥散性 \\
骨架式 & 碳骨架怎么连 & 碳链连接、官能团位置 & 碳和氢、真实键角、三维 \\
楔线式 & 谁朝前谁朝后 & 关键部位的立体指向 & 精确键角、旋转中的其他姿态 \\
球棍模型 & 骨架什么几何 & 键角、整体构型 & 真实比例、电子云 \\
空间填充模型 & 分子占多大地方 & 真实大小、表面形状 & 键、内部原子、连接方式 \\
\bottomrule
\end{tabular}
\caption{五种分子模型各自的信息交易。图中的颜色、大小、棍和线都是人为约定，不代表原子的真实外观。}
\end{table}

这张表值得记住的结论只有一句：\textbf{没有“最好”的模型，只有“适合当前问题”的模型}。第 38 章拆解有机反应时用骨架式，第 42 章谈手性时用楔线式，第 49 章讲酶的口袋时用空间填充模型——回头看，这本书一直在根据要回答的问题切换画像。你以后读任何科学文章，第一件事就是认出它用的是哪张画像，第二件事是问自己：这个问题，用这张画像回答合适吗？

\begin{shiyan}
\textbf{亲手感受“形状”：软糖分子模型}

材料：几种颜色的软糖（或橡皮泥搓的小球）、一盒牙签。

步骤：1. 搭水分子：一颗“氧”球、两颗“氢”球，牙签连接，把两根牙签掰成钝角（约 $105^{\circ}$）——别把两个氢放在氧的正对面；2. 搭二氧化碳：一颗“碳”、两颗“氧”，这次让两个氧在碳的两侧排成一条直线；3. 搭甲烷：一颗“碳”、四颗“氢”，试着让四根牙签两两之间的夹角尽量相等——你会发现它们自然张开成一个立体的架，怎么压都压不成平面上的十字。

观察点：甲烷的四根“键”为什么不肯待在同一个平面里？（提示：本节前面的桥梁框。）把水和二氧化碳并排摆好，猜一猜哪个分子有“正极”和“负极”之分，哪个没有——第 20、22 章会揭晓答案。

安全提示：牙签尖锐，别对着自己和别人；用过的软糖沾了灰尘和细菌，\textbf{不要吃掉}。
\end{shiyan}

\begin{wuqu}
“分子模型里哪种最像真的？我要学最真的那种。”——这个问题的前提就错了。空间填充模型最像分子的“轮廓”，可它对键和内部结构一无所知；路易斯结构最不像分子，却能告诉你电子的底细。这就像问“世界地图和地铁图哪个更真”：世界地图忠于形状和距离，地铁图忠于换乘关系——拿着世界地图进地铁站，照样迷路。模型的“真”从来不是像不像照片，而是\textbf{它保留的信息，是不是你要回答的那个问题所需要的}。同理，“蓝色的球就是蓝色的原子”也是把约定当成了事实：原子没有颜色，蓝色只是制图师发给氮的工作服。
\end{wuqu}

\begin{zhengju}
分子模型不只是教学挂图，它们曾经是被嘲笑的预言。1874 年，荷兰青年化学家范特霍夫还是个无名之辈，他发表了一本小册子，主张碳原子的四根键指向正四面体的四个顶点。证据是间接的：如果四根键躺在同一平面，二氯甲烷 \chem{CH_2Cl_2} 这类分子应该存在两种排列不同的版本，可化学家从来只找到一种——四面体模型恰好解释了“为什么只有一种”。当时的大化学家科尔贝公开嘲讽这是“坐在纸面上的空想”。

但这个纸面模型继续作出可被推翻的预言：连有四个不同基团的碳，应该存在左右手两种分子。此后几十年，乳酸等分子的旋光性（让偏振光旋转的本领，第 42 章）逐一兑现了这个预言。又过了近四十年，1912 年劳厄与布拉格父子开创 X 射线晶体学，人们第一次直接“看见”晶体中原子的真实排列——键角、键距，与当年纸面模型的预言吻合。范特霍夫因此获得了 1901 年首届诺贝尔化学奖。这个故事里最值得记住的是：一个好的分子模型不是图画，而是一台预言机——它敢说出还没人看见的东西，然后接受证据的审判。
\end{zhengju}

\section{这五种模型也有边界}

最后说说它们共同失效的地方。遇到电子离域的分子（苯环，第 41 章），任何“一根棍一根键”的画面都会误导你——那六对电子不属于任何两个原子之间；遇到金属（第 21 章），根本没有明确的“分子”可画；遇到蛋白质这样由成千上万个原子组成的庞然大物，连空间填充模型都会变成一团无法辨认的“菜花”，结构生物学只好再发明飘带图之类的第六种、第七种表示法（第 48 章）。

所以，每当你看见一张分子图，养成习惯先问两句话：它想让我看见什么？它为了不让我分心，藏起了什么？会这样追问的人，才算真正看懂了分子模型——也看懂了这本书接下来要给你看的每一张图。

\section*{练一练}

\begin{sikaoti}
\item 葡萄糖的分子式是 \chem{C_6H_{12}O_6}。同一个分子，分别用路易斯结构和空间填充模型表示，哪一种能让你数出“有几个羟基（$-\chem{OH}$）”？哪一种能让你判断“它能不能挤过细胞膜上的某个通道”？这个对比说明了什么？
\item 一张骨架式上有一个没有任何字母标记的折点。这个折点代表什么原子？它身上连着几个氢？你用了哪条规则才答得出来？
\item 两个互为镜像的乳酸分子，用普通骨架式画得出区别吗？必须换成哪种模型？为什么这个区别在药物里可能性命攸关？（提示：回想第 42 章。）
\item 球棍模型把“棍”画得又细又长，严重夸大了原子间距。假如把碳原子画成直径 2 厘米的球，按真实比例，与它成键的氢原子球（直径约 1.5 厘米）的中心应该离碳球中心多远？（提示：两球要深深嵌进彼此。）这个比例，五种模型里哪一种画对了？
\item 找一篇介绍药物或营养成分的科普文章，认出里面那张分子图属于五种模型中的哪一种，然后回答：作者想让你看见什么？这张图藏起了什么？换哪种模型能把藏起的信息补回来？
\end{sikaoti}
